\documentclass[main.tex]{subfiles}


\begin{document}

%from pagenumber to up
% \phantomsection 
% \hypertarget{toc}{}

 

 

 
   

\section{Сохранение и изменение энергии}


\textbf{Полная механическая энергия} ($E$ [Дж])~--- это сумма кинетической и потенциальной энергии тела\footnote{Для системы тел~--- сумма кинетических и потенциальных энергий тел системы.}:
\begin{equation}
    E=E_\text{к}+E_\text{п}.
\end{equation}
%
\begin{law}
\textbf{Закон сохранения полной механической энергии.} Если на тело (в системе тел) не действуют силы трения и внешние силы, то полная механическая энергия тела (системы) сохраняется\footnote{Вообще говоря, этот закон справедлив в случае равенства нулю суммы работ сил трения и внешних сил: $A_\text{тр}+A_\text{внеш}=0$.}:
\begin{equation}\label{45COE_e1}
    E_1=E_2=\ldots,
\end{equation}
где~$E_1$ и $E_2$~--- полные механические энергии в первом и втором состояниях.
\end{law}

Пусть, например, вначале шар покоится на горке, и сжатая пружина с приставленным к ней бруском неподвижна (рис.\,\ref{fig:p59}).

\begin{figure}[H]
    \centering

    \begin{tikzpicture}[font=\small,            line cap=round,                             >={Stealth[inset=0pt,round,length=3mm, angle'=20]},vec/.style={>={Stealth[inset=0pt,round,length=3.5mm, angle'=20]}}]


% \draw[lightgray,semitransparent,line width=.3cm,yshift=-.15cm,line cap=butt] (0,1.5) to[out=0,in=180] (4,0) -- (6,0); %не сходятся края толстой и тонкой
% \begin{scope}
%     \clip[] (0,1.5) to[out=0,in=180] (4,0) -- (6,0) --++ (0,-.3) --++ (-6,0) -- cycle; 
%     \draw[double=lightgray!50!white,lightgray!50!white,line cap=butt,double,double distance between line centers=.6cm] (0,1.5) to[out=0,in=180] (4,0) -- (6,0); %не сходятся края толстой и тонкой
% \end{scope}

\filldraw[lightgray,semitransparent] (0,1.5) to[out=0,in=180] (4,0) -- (6,0) --++ (0,-.3) --++ (-6,0) -- cycle;  
\draw (0,1.5) to[out=0,in=180] (4,0) -- (6,0);

%%%%%%%%%%%%%%%%%%%%%%%%%%%%%%
%%%%left pic
%solid
\shade[ball color=lightgray] (0,1.8) coordinate (C) circle (.3);

% h
\draw[densely dashed,gray] (0,1.5) --++ (5,0) coordinate (h1);
\draw[<->,gray] ([xshift=-.2cm]h1) --++ (0,-1.5) node[midway,right,black] {$h$};



%%%%%%%%%%%%%%%%%%%%%%%%%%%%%%%%
%%%%%%%%%%%%%%%%%%%%%%%%%%%%%%%%
%%%%%%%%%%%%%%%%%%%%%%%%%%%%%%%%
% right pic
\begin{scope}[xshift=8cm]
\filldraw[lightgray,semitransparent] (0,1.5) -- (0,0) -- (6,0) --++ (0,-.3) --++ (-6.3,0) --++ (0,1.8) -- cycle;
\draw[] (0,1.5) -- (0,0) --++ (6,0);

\filldraw[lightgray,draw=black] (1.4,0) rectangle++ (1,.5);

%% elast

% before
\draw (0,0)
foreach \x in {.1,.3,...,1.4}
{ -- ({\x},.5) -- ({\x+.1},0) }
;

% % after
% \draw[red] (0,.5) -- (0,0)
% foreach \x in {.2,.6,...,3}
% { -- ({\x},.5) -- ({\x+.2},0) }
% -- (2.8,.5)
% ;

% x
\draw[densely dashed,gray] (1.4,.5) --++ (0,1) coordinate (x1);
\draw[densely dashed,gray] (2.8,0) --++ (0,1.5);
\draw[<->,gray] ([yshift=-.2cm]x1) --++ (1.4,0) node[midway,above,black] {$x$};



\end{scope}

    \end{tikzpicture}
\caption{Начальные положения шара и сжатой пружины с бруском}
\label{fig:p59}
\end{figure}

Покоящийся шар (рис.\,\ref{fig:p59}, слева) массы~$m_\text{ш}$ на высоте~$h$ обладает в начальном положении полной механической энергией $E_\text{ш1}=m_\text{ш}gh$. Неподвижная система тел <<пружина"=брусок>> (рис.\,\ref{fig:p59}, справа) с деформацией~$x$ легкой пружины жесткости~$k$ обладает вначале полной механической энергией $E_\text{пб1}=\dfrac{kx^2}{2}$.

% (рис.\,\ref{fig:p59}).

% \begin{figure}[H]
%     \centering

%     \begin{tikzpicture}[font=\small,            line cap=round,                             >={Stealth[inset=0pt,round,length=3mm, angle'=20]},vec/.style={>={Stealth[inset=0pt,round,length=3.5mm, angle'=20]}}]


% % \draw[lightgray,semitransparent,line width=.3cm,yshift=-.15cm,line cap=butt] (0,1.5) to[out=0,in=180] (4,0) -- (6,0); %не сходятся края толстой и тонкой
% % \begin{scope}
% %     \clip[] (0,1.5) to[out=0,in=180] (4,0) -- (6,0) --++ (0,-.3) --++ (-6,0) -- cycle; 
% %     \draw[double=lightgray!50!white,lightgray!50!white,line cap=butt,double,double distance between line centers=.6cm] (0,1.5) to[out=0,in=180] (4,0) -- (6,0); %не сходятся края толстой и тонкой
% % \end{scope}

% \filldraw[lightgray,semitransparent] (0,1.5) to[out=0,in=180] (4,0) -- (6,0) --++ (0,-.3) --++ (-6,0) -- cycle;  
% \draw (0,1.5) to[out=0,in=180] (4,0) -- (6,0);

% %%%%%%%%%%%%%%%%%%%%%%%%%%%%%%
% %%%%left pic
% %solid
% \shade[ball color=lightgray] (5,.3) coordinate (C) circle (.3);

% % % h
% % \draw[densely dashed,gray] (0,1.5) --++ (5,0) coordinate (h1);
% % \draw[<->,gray] ([xshift=-.2cm]h1) --++ (0,-1.5) node[midway,right,black] {$h$};

% % v
% \draw[->,NavyBlue,thick,vec] (C) coordinate (F)--++ (0:1) node[black,above] {$\vec v_\text{ш}$};
% \node[] at (F) [circle,fill=NavyBlue,inner sep=0pt,minimum size=1mm,label=] {};



% %%%%%%%%%%%%%%%%%%%%%%%%%%%%%%%%
% %%%%%%%%%%%%%%%%%%%%%%%%%%%%%%%%
% %%%%%%%%%%%%%%%%%%%%%%%%%%%%%%%%
% % right pic
% \begin{scope}[xshift=8cm]
% \filldraw[lightgray,semitransparent] (0,1.5) -- (0,0) -- (6,0) --++ (0,-.3) --++ (-6.3,0) --++ (0,1.8) -- cycle;
% \draw[] (0,1.5) -- (0,0) --++ (6,0);

% \filldraw[lightgray,draw=black] (2.8,0) rectangle++ (1,.5);

% %% elast

% % % before
% % \draw (0,.5) -- (0,0)
% % foreach \x in {.1,.3,...,1.4}
% % { -- ({\x},.5) -- ({\x+.1},0) }
% % -- (1.4,.5)
% % ;

% % after
% \draw[] (0,.5) -- (0,0)
% foreach \x in {.2,.6,...,3}
% { -- ({\x},.5) -- ({\x+.2},0) }
% -- (2.8,.5)
% ;

% % % x
% % \draw[densely dashed,gray] (1.4,.5) --++ (0,1) coordinate (x1);
% % \draw[densely dashed,gray] (2.8,0) --++ (0,1.5);
% % \draw[<->,gray] ([yshift=-.2cm]x1) --++ (1.4,0) node[midway,above,black] {$x$};

% % v
% \draw[->,NavyBlue,thick,vec] (3.3,.25) coordinate (F)--++ (0:1) node[black,above] {$\vec v_\text{б}$};
% \node[] at (F) [circle,fill=NavyBlue,inner sep=0pt,minimum size=1mm,label=] {};


% \end{scope}

%     \end{tikzpicture}
% \caption{Конечные положения тел}
% \label{fig:p59}
% \end{figure}

Пусть рассматриваемые тела изменили свое положение.
Соскользнув без трения с горки на горизонтальную поверхность, шар приобрел скорость~$v_\text{ш}$, и его полная механическая энергия стала равна $E_\text{ш2}=\dfrac{m_\text{ш}^{\vphantom 2}v_\text{ш}^2}{2}$. Пружина возвращаясь в недеформированное состояние, привела в движение соприкасающийся с ней на гладкой поверхности брусок массы~$m_\text{б}$, который приобрел скорость~$v_\text{б}$; тогда полная механическая энергия этой системы равна $E_\text{пб2}=\dfrac{m_\text{б}^{\vphantom 2}v_\text{б}^2}{2}$. 

% Опыт показывает, что в рассмотренном примере полная механические энергии тела (системы тел) в начальном и конечном положениях равны.

Так, для рассмотренного процесса с шаром  соотношение~\eqref{45COE_e1} дает: $E_\text{ш1}\hm=E_\text{ш2}$ или $m_\text{ш}gh=\dfrac{m_\text{ш}^{\vphantom 2}v_\text{ш}^2}{2}$ (аналогично для системы <<пружина"=брусок>>).

Если же на тело (в системе тел) действуют силы трения и$/$или внешние силы, то справедлив следующий закон.
%
\begin{law}
\textbf{Закон изменения полной механической энергии.} Работа \emph{непотенциальных сил}\footnote{Под непотенциальными силами здесь можно понимать силы трения и внешние силы, тогда работа непотенциальных сил равна $A_\text{н.\,с}=A_\text{тр}+A_\text{внеш}$.}, действующих на тело (в системе тел), равна изменению полной механической энергии тела (системы):
\begin{equation}
    A_\text{н.\,с}=\Delta E.
\end{equation}
\end{law}

% Так для рассмотренных процессов с шаром и системой <<пружина--брусок>> соотношение~\eqref{45COE_e1} дает: $E_\text{ш1}=E_\text{ш2}$ или 
















 
\end{document}